% Provenance (never printed): the demography strand of ../economics_series.md; urbanisation
% added from the geography programme (one paragraph). Brief: sources/book1/BRIEFS.md, g6-04.
% Chart data: figdata/grade-6/04-health-and-population (Gapminder lex and pop before 1950;
% WDI SP.DYN.LE00.IN, SP.POP.TOTL, SH.DYN.MORT, NY.GDP.PCAP.PP.KD, SP.DYN.TFRT.IN,
% SP.URB.TOTL.IN.ZS, SP.POP.DPND.OL).
\chapter{Longer Lives, Smaller Families}\label{ch:g6:health-and-population}

The world's population rose from one billion to eight billion in two
centuries. One might guess that families had more children. The opposite
happened: families now have far fewer children than in 1800. What changed is
that most children now live to grow up, and most adults live to grow old. This
chapter measures both changes, and then follows what came next --- smaller
families and older populations.

\begin{omfigure}
\includegraphics[width=0.62\linewidth]{images/book1/ai/g6-health-and-population-family.jpg}

{\small Three generations at one table. In 1800, few children knew their
grandparents for long; today most do.}
\end{omfigure}

\begin{recall}
A \omterm{def:g6:reading-economic-charts:per-person}{per-person value} and a share were defined in
\cref{ch:g6:reading-economic-charts}, with the \omterm{def:g6:reading-economic-charts:times-ten-scale}{times-ten scale}; \omterm{def:g6:living-standards:income-per-person}{income per
person} in \cref{ch:g6:living-standards}.
\end{recall}

\section{Life expectancy at birth}

\begin{definition}[Life expectancy at birth]\label{def:g6:health-and-population:life-expectancy}
The \emph{life expectancy at birth}\index{life expectancy at birth} of a
country in a year is the average number of years that a newborn baby would
live if, all through its life, people of each age kept dying at the rates
observed in that country in that year.
\end{definition}

\begin{measure}[Life expectancy at birth]
Statisticians count, for each age, how many people of that age died during the
year and how many were alive. From these death rates they build a \emph{life
table}: out of 100\,000 imaginary babies, how many would die before age one,
between one and two, and so on, up to the oldest ages. Life expectancy is the
average age at death in that table. Two things follow. It is an average:
many deaths of young children pull it down a great deal, even if the adults
live long. And it is a snapshot of one year: in a year of war or epidemic it
falls sharply, although the babies born that year will not live through the
war or the epidemic all their lives.
\end{measure}

\begin{omchart}{Life expectancy at birth, 1800--2024, in years, in six world regions
and the world. Before 1950 the values are Gapminder's estimates; from 1960 they
are the World Bank's. The years 1950--1959 are left empty.}{Gapminder
Foundation (life expectancy and population), 1800--1949; World Bank, World
Development Indicators, SP.DYN.LE00.IN and SP.POP.TOTL, 1960--2024; regional
values weighted by population}
\begin{tikzpicture}
\begin{axis}[omtime, xmin=1800, xmax=2024, ymin=0, ymax=85,
  xtick={1800,1850,1900,1950,2000}, xlabel={Year}, ylabel={Years},
  height=0.6\linewidth, legend style={at={(0.5,-0.2)}, anchor=north, legend columns=3},
  ]
\addplot[omDef, thick] table[x=year, y=america]
  {figdata/grade-6/04-health-and-population/out/life_expectancy_early.csv};
\addplot[omMeth, thick] table[x=year, y=east_asia_pacific]
  {figdata/grade-6/04-health-and-population/out/life_expectancy_early.csv};
\addplot[omProp, thick] table[x=year, y=europe_central_asia]
  {figdata/grade-6/04-health-and-population/out/life_expectancy_early.csv};
\addplot[omDebate, thick] table[x=year, y=middle_east_north_africa]
  {figdata/grade-6/04-health-and-population/out/life_expectancy_early.csv};
\addplot[omExo, thick] table[x=year, y=south_asia]
  {figdata/grade-6/04-health-and-population/out/life_expectancy_early.csv};
\addplot[omThm, thick] table[x=year, y=sub_saharan_africa]
  {figdata/grade-6/04-health-and-population/out/life_expectancy_early.csv};
\addplot[black, thick, densely dashed] table[x=year, y=world]
  {figdata/grade-6/04-health-and-population/out/life_expectancy_early.csv};
\legend{America, East Asia and Pacific, Europe and Central Asia,
  Middle East and North Africa, South Asia, Sub-Saharan Africa, World}
\addplot[omDef, thick, forget plot] table[x=year, y=america]
  {figdata/grade-6/04-health-and-population/out/life_expectancy_late.csv};
\addplot[omMeth, thick, forget plot] table[x=year, y=east_asia_pacific]
  {figdata/grade-6/04-health-and-population/out/life_expectancy_late.csv};
\addplot[omProp, thick, forget plot] table[x=year, y=europe_central_asia]
  {figdata/grade-6/04-health-and-population/out/life_expectancy_late.csv};
\addplot[omDebate, thick, forget plot] table[x=year, y=middle_east_north_africa]
  {figdata/grade-6/04-health-and-population/out/life_expectancy_late.csv};
\addplot[omExo, thick, forget plot] table[x=year, y=south_asia]
  {figdata/grade-6/04-health-and-population/out/life_expectancy_late.csv};
\addplot[omThm, thick, forget plot] table[x=year, y=sub_saharan_africa]
  {figdata/grade-6/04-health-and-population/out/life_expectancy_late.csv};
\addplot[black, thick, densely dashed, forget plot] table[x=year, y=world]
  {figdata/grade-6/04-health-and-population/out/life_expectancy_late.csv};
\end{axis}
\end{tikzpicture}
\end{omchart}

% figdata: life_expectancy_early (world 1800 30.6, 1900 33.6, 1917 35.7, 1918 24.3, 1949 50.7;
% 1800 regions 25.7 to 33.9; south_asia 1918 9.8)
% figdata: life_expectancy_late (world 1960 50.9, 2019 72.9, 2021 71.2, 2024 73.5;
% 2024: europe_central_asia 78.6, sub_saharan_africa 62.8)
In 1800, \omterm{def:g6:health-and-population:life-expectancy}{life expectancy at birth} was between 26 and 34 years in every region,
and $30.6$ years for the world. It barely moved for a century: $33.6$ years in
1900. Then it rose almost everywhere: $50.9$ years for the world in 1960 and
$73.5$ years in 2024. The highest regional value in 2024 was $78.6$ years
(Europe and Central Asia), the lowest $62.8$ years (Sub-Saharan Africa) --- far
above any region in 1800.

The deep notch in 1918 is the influenza pandemic of 1918--1919, which killed
tens of millions of people in a few months. The world's life expectancy fell
from $35.7$ years in 1917 to $24.3$ in 1918, then came back the next year.
The estimates for South Asia fall to about 10 years in 1918. A smaller dip in
2020 and 2021 is the COVID-19 pandemic: from $72.9$ years in 2019 to $71.2$ in
2021.

\section{Child mortality}

\begin{definition}[Child mortality rate]\label{def:g6:health-and-population:child-mortality}
The \emph{child mortality rate}\index{child mortality rate} is the number of
children who die before their fifth birthday, for every 1\,000 children born
alive.
\end{definition}

A \omterm{def:g6:health-and-population:child-mortality}{child mortality rate} of 100 means that one child in ten dies before the age
of five. Around 1800, the estimates for every country are several hundred:
% figdata: facts (cm_1800_min 322, cm_1800_max 540)
from 322 to 540 deaths per 1\,000 births, so that between a third and a half of
all children died before five.
This is the main reason why \omterm{def:g6:health-and-population:life-expectancy}{life expectancy at birth} was so low.

\begin{omchart}{Child mortality and income per person in 2024, one dot per
country: deaths before age five per 1\,000 births, against income per person in
dollars of 2021 on a times-ten scale.}{World Bank, World Development Indicators,
SH.DYN.MORT and NY.GDP.PCAP.PP.KD, 2024}
\begin{tikzpicture}
\begin{semilogxaxis}[omincomex, xmin=700, xmax=200000, ymin=0, ymax=120,
  xlabel={Income per person (dollars of 2021, times-ten scale)},
  ylabel={Deaths before age 5 per 1\,000 births}]
\addplot table[x=income, y=mortality]
  {figdata/grade-6/04-health-and-population/out/child_mortality_income.csv};
\end{semilogxaxis}
\end{tikzpicture}
\end{omchart}

% figdata: child_mortality_income (median under 3000: 62; above 30000: 4.3)
% figdata: child_mortality_pairs (GNQ 14998 68; LKA 13793 5.9; DOM 24230 30.6; MKD 24203 2.8)
The dots fall from left to right: the richer a country, the fewer of its children
die. In the countries with less than 3\,000 dollars per person, the middle
value is about 62 deaths per 1\,000 births; in those with more than 30\,000
dollars, about 4. But at the same income the dots spread widely. Equatorial
Guinea, with $14\,998$ dollars per person, had 68 deaths per 1\,000 births; Sri
Lanka, with $13\,793$ dollars, had $5.9$. Income helps a great deal, and it is
not the whole story: clean water, vaccines, clinics and the schooling of mothers
matter too.

\section{Births per woman}

\begin{definition}[Births per woman]\label{def:g6:health-and-population:births-per-woman}
The number of \emph{births per woman}\index{births per woman (fertility rate)}
in a country in a year, also called its \emph{fertility rate}, is the number
of children a woman would have in her life if, at each age, she had children
at the rates observed among women of that age in that year.
\end{definition}

\begin{measure}[Births per woman]
For each age from 15 to 49, the statistical office divides the births to
mothers of that age by the number of women of that age. Adding up these rates
over all ages gives the number of \omterm{def:g6:health-and-population:births-per-woman}{births per woman}. A value of about $2.1$ is
the \emph{replacement level}: each generation of parents is just replaced by its
children (a little more than 2, because some children die before they have
children of their own).
\end{measure}

\begin{omchart}{Births per woman by world region and in the world, 1960--2024.}{World
Bank, World Development Indicators, SP.DYN.TFRT.IN, World Bank regions}
\begin{tikzpicture}
\begin{axis}[omtime, xmin=1960, xmax=2024, ymin=0, ymax=8,
  xtick={1960,1970,1980,1990,2000,2010,2020}, xlabel={Year}, ylabel={Births per woman},
  height=0.6\linewidth, legend style={at={(0.5,-0.2)}, anchor=north, legend columns=2},
  cycle list={{omDef, thick}, {omMeth, thick}, {omProp, thick}, {omDebate, thick},
    {omDef, thick, densely dotted}, {omExo, thick}, {omThm, thick}, {black, thick, densely dashed}}]
\pgfplotsforeachungrouped \col in {eas, ecs, lcn, mea, nac, sas, ssf, wld} {
  \addplot table[x=year, y=\col]
    {figdata/grade-6/04-health-and-population/out/fertility_regions.csv};
}
\legend{East Asia and Pacific, Europe and Central Asia, Latin America and Caribbean,
  {Middle East, North Africa, Afghanistan, Pakistan}, North America, South Asia,
  Sub-Saharan Africa, World}
\end{axis}
\end{tikzpicture}
\end{omchart}

% figdata: fertility_regions (wld 1960 4.69, 2024 2.19, first below 3 in 1993; ssf 1960 6.57, 2024 4.26;
% eas 2024 1.34; sas 1960 6.00, 2024 1.98)
In 1960, a woman had on average $4.69$ children in the world as a whole. In
2024, $2.19$. The fall was largest in East Asia, now the region with the fewest
births ($1.34$ per woman), and in South Asia, from $6.00$ to $1.98$. Sub-Saharan
Africa's rate has fallen too, from $6.57$ to $4.26$, but it is still the
highest. In most regions, women now have fewer children than the replacement
level.

\section{Why population grew fast, and why its growth is slowing}

A population is like a tank of water. Births flow in; deaths flow out. People
also move from one country to another, but for the world as a whole they cancel.

\begin{omfigure}
\begin{tikzpicture}[econ, x=1cm, y=1cm]
  % invented town: 10 000 people, 300 births and 100 deaths a year
  \draw[thick] (0,0) -- (0,3) (4,3) -- (4,0) -- (0,0);
  \fill[omDef!20] (0,0) rectangle (4,2.2);
  \node[font=\footnotesize, align=center] at (2,1.1) {population\\10\,000};
  \draw[omMeth, -{Stealth[length=3mm]}, line width=2.5pt] (-1.6,3.6) to[out=0, in=120] (0.8,2.4);
  \node[font=\footnotesize, text=omMeth, align=center] at (-1.8,4.0) {births: 300 a year};
  \draw[omThm, -{Stealth[length=3mm]}, line width=1.2pt] (4,0.4) -- (5.6,0.4);
  \node[font=\footnotesize, text=omThm, anchor=west] at (5.7,0.4) {deaths: 100 a year};
  \node[font=\footnotesize, align=left, anchor=west] at (5.7,2.2) {one year later:\\$10\,000 + 300 - 100 = 10\,200$};
\end{tikzpicture}

{\small A population as a tank (an invented town). Its change in a year is the
births minus the deaths, plus the people who arrive minus those who leave.}
\end{omfigure}

Before 1800, births and deaths were both high, about equal on average, and the
tank barely rose. Then deaths fell, first among children: clean water, better
food, vaccines. Births stayed high for a generation or more, because families
were used to having many children, many of whom used to die. For those decades
the tank filled fast: this is why the world's population rose from 2.5 billion
in 1950 to 6 billion in 2000. Then births fell too. Parents whose children
survive have fewer of them, and parents who move to towns, where children go to
school instead of working on a farm, have fewer still. With fewer births, the
population grows more slowly.

% figdata: facts (urban_1960 34.2, urban_2025 57.8)
The move to towns is itself one of the large changes of these years. In 1960,
$34.2\,\%$ of the world's people lived in towns and cities; in 2025, $57.8\,\%$.

\section{An ageing population}

When fewer children are born and adults live longer, the share of older people
in the population rises. Economists watch one ratio closely, because people of
working age produce most of what everyone, young and old, consumes.

\begin{definition}[Old-age dependency ratio]\label{def:g6:health-and-population:old-age-dependency}
The \emph{old-age dependency ratio}\index{old-age dependency ratio} of a
population is the number of people aged 65 and over for every 100 people aged
15 to 64.
\end{definition}

\begin{omfigure}
\begin{tikzpicture}[econ, x=1cm, y=1cm]
  % invented towns of 10 000; bars: thousands of people by age group (scale 0.5 cm per 1 000)
  \foreach \i/\g/\a/\b in {0/{0--14}/4.0/1.5, 1/{15--39}/3.5/2.8, 2/{40--64}/2.0/3.2,
      3/{65--79}/0.4/1.9, 4/{80+}/0.1/0.6} {
    \fill[omDef!50] (0,{\i*0.6}) rectangle ({\a*0.5},{\i*0.6+0.45});
    \node[left, font=\scriptsize] at (-0.1,{\i*0.6+0.22}) {\g};
    \fill[omThm!50] (6.2,{\i*0.6}) rectangle ({6.2+\b*0.5},{\i*0.6+0.45});
    \node[left, font=\scriptsize] at (6.1,{\i*0.6+0.22}) {\g};
  }
  \node[font=\footnotesize, align=center] at (1.2,-0.6) {young town:\\65+: 500;\quad 15--64: 5\,500};
  \node[font=\footnotesize, align=center] at (7.4,-0.6) {ageing town:\\65+: 2\,500;\quad 15--64: 6\,000};
  \node[font=\scriptsize] at (1.2,3.25) {age group};
  \node[font=\scriptsize] at (7.4,3.25) {age group};
\end{tikzpicture}

{\small Two invented towns of 10\,000 people, by age group: each bar is the
number of people in the group. The young town has 9 people aged 65 and over for
every 100 of working age ($500 \div 5\,500 \times 100 \approx 9$); the ageing
town has 42 ($2\,500 \div 6\,000 \times 100 \approx 42$).}
\end{omfigure}

% figdata: facts (old_dep_1960 8.6, old_dep_2024 15.7)
For the world as a whole, the \omterm{def:g6:health-and-population:old-age-dependency}{old-age dependency ratio} rose from $8.6$ in 1960
to $15.7$ in 2024: from about one older person for every twelve of working age
to about one for every six. In countries where births fell early, it is much
higher. An ageing population is a success --- people live longer --- and it
raises new questions: who will care for the old, and how will pensions be paid?
Later books of this course return to them.

\section{Exercises}

\begin{exercise}[$\star$]\label{exo:g6:health-and-population:1}
Use the chart of life expectancy. About what was the world's \omterm{def:g6:health-and-population:life-expectancy}{life expectancy at
birth} in 1900, and in 2024?
\end{exercise}

\begin{exercise}[$\star$]\label{exo:g6:health-and-population:2}
Use the chart of \omterm{def:g6:health-and-population:births-per-woman}{births per woman}. Which region had the fewest \omterm{def:g6:health-and-population:births-per-woman}{births per woman}
in 2024? Which had the most?
\end{exercise}

\begin{exercise}[$\star$]\label{exo:g6:health-and-population:3}
In an invented country, 37 children die before age five for every 1\,000 born.
(a) What percentage of children is that? (b) Out of 50\,000 children born in a
year, about how many will die before five?
\end{exercise}

\begin{exercise}[$\star$]\label{exo:g6:health-and-population:4}
Use the chart of child mortality and income. (a) Roughly how many children per
1\,000 die before five in the countries near 2\,000 dollars per person? (b) In
the countries near 50\,000 dollars?
\end{exercise}

\begin{exercise}[$\star\star$]\label{exo:g6:health-and-population:5}
In an invented village, ten people died one year, at these ages: 0, 0, 0, 2, 4,
40, 55, 60, 68 and 71. (a) Compute the average age at death. (b) How many of the
ten died at about that age? (c) Compute the average age at death of those who
reached the age of five. What does this tell you about ``a life expectancy of
30 years''?
\end{exercise}

\begin{exercise}[$\star\star$]\label{exo:g6:health-and-population:6}
An invented town has 1\,200 people aged 65 and over, and 4\,800 people aged 15 to
64. Compute its \omterm{def:g6:health-and-population:old-age-dependency}{old-age dependency ratio}.
\end{exercise}

\begin{exercise}[$\star\star$]\label{exo:g6:health-and-population:7}
Use the chart of \omterm{def:g6:health-and-population:births-per-woman}{births per woman}. In about which year did the world's line fall
below 3 \omterm{def:g6:health-and-population:births-per-woman}{births per woman}?
\end{exercise}

\begin{exercise}[$\star\star$]\label{exo:g6:health-and-population:8}
Someone says: ``A richer country always has a lower \omterm{def:g6:health-and-population:child-mortality}{child mortality rate} than a
poorer one.'' (a) What two dots on the chart of child mortality would contradict
this sentence? (b) Does the chart contain such dots? (c) Write a sentence that
the chart does support.
\end{exercise}

\begin{exercise}[$\star\star$]\label{exo:g6:health-and-population:9}
An invented town has 5\,000 people. In the first year there are 150 births and 60
deaths; in the second, 140 births and 65 deaths; in the third, 130 births and 70
deaths. Nobody arrives or leaves. (a) Compute the population at the end of each
year. (b) Is the population still growing? Is its growth speeding up or slowing
down?
\end{exercise}

\begin{exercise}[$\star\star\star$]\label{exo:g6:health-and-population:10}
% figdata: child_mortality_pairs (DOM 24230 30.6; MKD 24203 2.8)
In 2024, the Dominican Republic had an \omterm{def:g6:living-standards:income-per-person}{income per person} of $24\,230$ dollars and
$30.6$ deaths before age five per 1\,000 births; North Macedonia had $24\,203$
dollars and $2.8$ deaths. (a) By how many dollars do their incomes differ?
(b) How many times higher is the Dominican Republic's child mortality? Round to
one decimal. (c) Name two things, other than \omterm{def:g6:living-standards:income-per-person}{income per person}, that could
explain such a gap.
\end{exercise}

\begin{exercise}[$\star\star\star$]\label{exo:g6:health-and-population:11}
An invented town has 2\,000 people aged 65 and over and 6\,000 aged 15 to 64.
Every five years, the number aged 65 and over rises by 300 and the number aged 15
to 64 falls by 200. Compute the \omterm{def:g6:health-and-population:old-age-dependency}{old-age dependency ratio} now and after 5, 10, 15
and 20 years, to one decimal. Is it rising by the same amount every five years?
\end{exercise}

\begin{exercise}[$\star\star\star$]\label{exo:g6:health-and-population:12}
The chart of life expectancy gives one value per region and year. (a) Name two
groups of people, inside one country, whose life expectancies could be very
different. (b) What data would you need to compare them? (c) Can the regional
chart tell you whether the gap between them is widening?
\end{exercise}

\section{Problem: The Village Register}

\begin{problem}[{Weekend problem --- child mortality and average age at death from
a register, a village over thirty years, and the dependency ratio now and
later}]
\label{pb:g6:health-and-population:1}
The register of an invented village records 40 births in the year 1825 and 20
deaths, at these ages (in years):

\begin{center}
\begin{tabular}{ll}
\toprule
Children under five & 0, 0, 0, 0, 0, 1, 1, 2, 3, 4 \\
Older people        & 15, 28, 35, 42, 50, 58, 63, 67, 72, 81 \\
\bottomrule
\end{tabular}
\end{center}

\medskip
\noindent\textbf{Part I --- One year of the register.}
\begin{enumerate}
  \item How many children under five died in 1825?
  \item Use the deaths under five and the births of the year to estimate the
        \omterm{def:g6:health-and-population:child-mortality}{child mortality rate}, per 1\,000 births.
  \item Compute the average age at death of the 20 people, to one decimal.
  \item Compute the average age at death of the 10 people who had reached five.
        Explain the difference with question 3.
\end{enumerate}

\medskip
\noindent\textbf{Part II --- Thirty years.} The village had 800 people in 1825.
From 1825 to 1835 it had 40 births and 20 deaths every year; from 1835 to 1845,
40 births and 12 deaths every year; from 1845 to 1855, 30 births and 12 deaths
every year. Nobody arrived or left.
\begin{enumerate}[resume]
  \item What was the population in 1835?
  \item In 1845?
  \item In 1855?
  \item Compute the \omterm{def:g6:reading-economic-charts:rate-of-change}{rate of change} of the population from 1825 to 1855. In which
        decade did it grow fastest, and why?
\end{enumerate}

\medskip
\noindent\textbf{Part III --- The village today.} Today the village has 220
people aged 0 to 14, 900 aged 15 to 64 and 380 aged 65 and over. In twenty
years, the town hall expects 180, 800 and 480.
\begin{enumerate}[resume]
  \item Compute the \omterm{def:g6:health-and-population:old-age-dependency}{old-age dependency ratio} today, to one decimal.
  \item Compute it in twenty years.
  \item How many people of working age are there for each person aged 65 and
        over, today and in twenty years? Round to two decimals.
  \item Give one reason why the town hall cares about the answer to
        question 11.
\end{enumerate}
\end{problem}
